% GENERATED FILE. Edit canonical lessons, then run npm run generate:books.
% canonical-source-hash: fnv1a64:d6d3becb17ab0a2f
% canonical-lessons: ZH-C79-shangke, ZH-W79-ke-lesson, ZH-C79-xiake, ZH-C79-jiaoshi, ZH-C79-zuoye, ZH-C79-kaoshi, ZH-R79-review-75-79

\chapter{At School}
\label{ch:zh-at-school}
\hlchaptermodality{79}

\begin{chapteropening}
\textbf{By the end of this chapter:} \emph{I can talk about classes, the classroom, homework and exams.}

Complete the last lesson of chapter 79: I can talk about classes, the classroom, homework and exams.
\end{chapteropening}

\section[\texorpdfstring{\zh{上课} (shàngkè)}{上课 (shàngkè)}]{\zh{上课} (shàngkè) — to have class; to start class}

\label{lesson:ZH-C79-shangke}



\begin{quote}
Before the new one: say the Mandarin for age, then the Mandarin for in childhood.
\end{quote}

\subsection*{You'll want to know: \zh{上课}}
\textbf{\zh{上课}} — \emph{shàngkè} — \textquotedblleft{}to have class; to start class\textquotedblright{}.

\begin{quote}
\textbf{\zh{今天有汉语课我们一会儿上课}} \emph{jīntiān yǒu hànyǔ kè wǒmen yíhuìr shàngkè} — There's a Chinese class today; it starts in a moment.
\end{quote}

The new character, in pieces: \zh{课} — \zh{讠} on the left, and on the right guǒ, the fruit of shuǐguǒ: a box over the tree piece. The next lesson puts it in your hand.

\begin{cousinweb}
\textbf{\zh{上课}} (\emph{shàngkè}) is built from \zh{上} 'to go up to' and \zh{课} 'lesson', both inherited from Middle Chinese, \emph{dzyangX} and \emph{khwaH} (Guangyun, in Baxter's transcription). The dictionary glosses the second as 'a tax' and 'a test'; the sense lesson is not given there.
\end{cousinweb}

\begin{grammarlens}[title={What you've built}]
One more word about people.
\end{grammarlens}

\subsection*{Your turn}
\begin{itemize}
\raggedright
  \item \emph{Say it:} \emph{shàngkè}
  \item \emph{Say it:} \emph{shàngkè}, once more
  \item \emph{Say it:} \emph{jīntiān yǒu hànyǔ kè wǒmen yíhuìr shàngkè} — There's a Chinese class today; it starts in a moment.
  \item \emph{Recall:} say the Mandarin for age, then the Mandarin for in childhood, then say \emph{shàngkè} again
\end{itemize}

\begin{tcolorbox}[breakable,colback=teal!4,colframe=teal!35!black,title={Before you move on}]
What does \zh{上课} mean? (\textquotedblleft{}To have class; to start class\textquotedblright{}.) Say it once more.
\end{tcolorbox}
\section[\texorpdfstring{\zh{课} (kè)}{课 (kè)}]{\zh{课} — the speech piece, then guǒ}

\label{lesson:ZH-W79-ke-lesson}



\begin{quote}
The word first, then the character inside it.

\begin{itemize}
\raggedright
  \item \emph{Recall:} say \emph{shàngkè} — \textbf{R1}, one lesson back, before it decays
\end{itemize}

\begin{quote}
\zh{讠} \zh{说}
\end{quote}

\emph{Read it:} each one aloud, with nothing said first — the speech piece, \emph{shuō}
\end{quote}

\subsection*{Script — the speech piece, then guǒ}
\hlblockfigure{\detokenize{figures/ZH-W79-ke-lesson-filmstrip.pdf}}{How \zh{课} is written, stroke by stroke}

\begin{quote}
\zh{课}
\end{quote}

\textbf{Two pieces, left before right.}

\textbf{\zh{讠}, two strokes}, as in \zh{说} and \zh{话}.

\textbf{Guǒ, eight strokes}: a box like \zh{日}, its left side, the top that turns down, the middle bar and the bar that closes it; then a long horizontal, a vertical down through the box and the bar, a left-falling stroke and a stroke falling to the right.

2 + 8 = ten.

\subsection*{Your turn}
\begin{itemize}
\raggedright
  \item \emph{Write it:} \zh{课}, all ten, with the model in front of you
  \item \emph{Cover:} write it again from memory, and count the pieces, not the strokes
\end{itemize}

\begin{tcolorbox}[breakable,colback=teal!4,colframe=teal!35!black,title={Before you move on}]
How many pieces in \zh{课}? (\textbf{Two}.) Which characters you know start the same way? (\textbf{\zh{说} and \zh{话}}.) What runs down through the box? (\textbf{The vertical, the eighth stroke}.)

Source: \href{https://raw.githubusercontent.com/chanind/hanzi-writer-data/68d10a4b21150cae5e1ebbd223eed289cf32d90c/data/\%E8\%AF\%BE.json}{Hanzi Writer Data \zh{课}, pinned stroke paths and medians}.
\end{tcolorbox}
\section[\texorpdfstring{\zh{下课} (xiàkè)}{下课 (xiàkè)}]{\zh{下课} (xiàkè) — to finish class}

\label{lesson:ZH-C79-xiake}



\begin{quote}
Before the new one: say the Mandarin for in childhood, then the Mandarin for to start class.
\end{quote}

\subsection*{You'll want to know: \zh{下课}}
\textbf{\zh{下课}} — \emph{xiàkè} — \textquotedblleft{}to finish class\textquotedblright{}.

\begin{quote}
\textbf{\zh{十二点前我们下课}} \emph{shí'èr diǎn qián wǒmen xiàkè} — We finish class before twelve.
\end{quote}

\begin{cousinweb}
\textbf{\zh{下课}} (\emph{xiàkè}) is built from \zh{下} 'to go down from' and \zh{课} 'lesson', both inherited from Middle Chinese, \emph{haeX} and \emph{khwaH} (Guangyun, in Baxter's transcription); lesson is not among the senses the dictionary gives the second.
\end{cousinweb}

\begin{grammarlens}[title={What you've built}]
One more word about people.
\end{grammarlens}

\subsection*{Your turn}
\begin{itemize}
\raggedright
  \item \emph{Say it:} \emph{xiàkè}
  \item \emph{Say it:} \emph{xiàkè}, once more
  \item \emph{Say it:} \emph{shí'èr diǎn qián wǒmen xiàkè} — We finish class before twelve.
  \item \emph{Recall:} say the Mandarin for in childhood, then the Mandarin for to start class, then say \emph{xiàkè} again
\end{itemize}

\begin{tcolorbox}[breakable,colback=teal!4,colframe=teal!35!black,title={Before you move on}]
What does \zh{下课} mean? (\textquotedblleft{}To finish class\textquotedblright{}.) Say it once more.
\end{tcolorbox}
\section[jiàoshì]{jiàoshì — classroom}

\label{lesson:ZH-C79-jiaoshi}



\begin{quote}
Before the new one: say the Mandarin for to start class, then the Mandarin for to finish class.
\end{quote}

\subsection*{You'll want to know: jiàoshì}
\emph{jiàoshì} — \textquotedblleft{}classroom\textquotedblright{}.

\begin{quote}
\emph{wǒmen de jiàoshì hěn dà} — Our classroom is big.
\end{quote}

Its characters come later; for now it is a word for the ear.

\begin{cousinweb}
\emph{jiàoshì} is built from jiào 'to teach' and shì 'room', both inherited from Middle Chinese, \emph{kaewH} and \emph{syit} (Guangyun, in Baxter's transcription).
\end{cousinweb}

\begin{grammarlens}[title={What you've built}]
One more word about people.
\end{grammarlens}

\subsection*{Your turn}
\begin{itemize}
\raggedright
  \item \emph{Say it:} \emph{jiàoshì}
  \item \emph{Say it:} \emph{jiàoshì}, once more
  \item \emph{Say it:} \emph{wǒmen de jiàoshì hěn dà} — Our classroom is big.
  \item \emph{Recall:} say the Mandarin for to start class, then the Mandarin for to finish class, then say \emph{jiàoshì} again
\end{itemize}

\begin{tcolorbox}[breakable,colback=teal!4,colframe=teal!35!black,title={Before you move on}]
What does jiàoshì mean? (\textquotedblleft{}Classroom\textquotedblright{}.) Say it once more.
\end{tcolorbox}
\section[zuòyè]{zuòyè — homework}

\label{lesson:ZH-C79-zuoye}



\begin{quote}
Before the new one: say the Mandarin for to finish class, then the Mandarin for classroom.
\end{quote}

\subsection*{You'll want to know: zuòyè}
\emph{zuòyè} — \textquotedblleft{}homework\textquotedblright{}.

\begin{quote}
\emph{jīntiān de zuòyè shì dú hànzì} — Today's homework is reading characters aloud.
\end{quote}

Its characters come later; for now it is a word for the ear.

\begin{cousinweb}
\emph{zuòyè} is built from zuò 'to do' and yè 'task', both inherited from Middle Chinese, \emph{tsak} and \emph{ngjaep} (Guangyun, in Baxter's transcription).
\end{cousinweb}

\begin{grammarlens}[title={What you've built}]
One more word about people.
\end{grammarlens}

\subsection*{Your turn}
\begin{itemize}
\raggedright
  \item \emph{Say it:} \emph{zuòyè}
  \item \emph{Say it:} \emph{zuòyè}, once more
  \item \emph{Say it:} \emph{jīntiān de zuòyè shì dú hànzì} — Today's homework is reading characters aloud.
  \item \emph{Recall:} say the Mandarin for to finish class, then the Mandarin for classroom, then say \emph{zuòyè} again
\end{itemize}

\begin{tcolorbox}[breakable,colback=teal!4,colframe=teal!35!black,title={Before you move on}]
What does zuòyè mean? (\textquotedblleft{}Homework\textquotedblright{}.) Say it once more.
\end{tcolorbox}
\section[kǎoshì]{kǎoshì — to take an exam}

\label{lesson:ZH-C79-kaoshi}



\begin{quote}
Before the new one: say the Mandarin for classroom, then the Mandarin for homework.
\end{quote}

\subsection*{You'll want to know: kǎoshì}
\emph{kǎoshì} — \textquotedblleft{}to take an exam\textquotedblright{}.

\begin{quote}
\emph{míngtiān kǎoshì jīntiān wǒ xuéxí} — The exam is tomorrow, so today I'm studying.
\end{quote}

Its characters come later; for now it is a word for the ear.

\begin{cousinweb}
\emph{kǎoshì} is built from kǎo and shì, both inherited from Middle Chinese, \emph{khawX} and \emph{syiH} (Guangyun, in Baxter's transcription). The dictionary glosses the first as 'to check' and the second as 'to use'; the sense 'exam' is not given there.
\end{cousinweb}

\begin{grammarlens}[title={What you've built}]
One more word about people.
\end{grammarlens}

\subsection*{Your turn}
\begin{itemize}
\raggedright
  \item \emph{Say it:} \emph{kǎoshì}
  \item \emph{Say it:} \emph{kǎoshì}, once more
  \item \emph{Say it:} \emph{míngtiān kǎoshì jīntiān wǒ xuéxí} — The exam is tomorrow, so today I'm studying.
  \item \emph{Recall:} say the Mandarin for classroom, then the Mandarin for homework, then say \emph{kǎoshì} again
\end{itemize}

\begin{tcolorbox}[breakable,colback=teal!4,colframe=teal!35!black,title={Before you move on}]
What does kǎoshì mean? (\textquotedblleft{}To take an exam\textquotedblright{}.) Say it once more.
\end{tcolorbox}
\section[(dialogue)]{Looking back: from Celebrating to At School}

\label{lesson:ZH-R79-review-75-79}



\begin{quote}
Say the Mandarin for joyful and birthday.
\end{quote}

\subsection*{Your turn}
\begin{itemize}
\raggedright
  \item \emph{Say it:} joyful, birthday, cake, present, to give as a present
  \item \emph{Say it:} a little while, see you tomorrow, to leave it till later, good night, bye-bye
  \item \emph{Say it:} don't mention it, to trouble somebody, to thank, extremely, really
  \item \emph{Say it:} how old, to grow up, ago, age, in childhood
  \item \emph{Say it:} to start class, to finish class, classroom, homework, to take an exam
  \item \emph{Read it:} \zh{乐} \zh{会} \zh{前} \zh{课} — each one aloud, then a word that holds it
\end{itemize}

\begin{tcolorbox}[breakable,colback=teal!4,colframe=teal!35!black,title={Before you move on}]
Joyful? (\textbf{\zh{快乐}}, \emph{kuàilè}.) Birthday? (\textbf{\zh{生日}}, \emph{shēngrì}.) Cake? (\emph{dàngāo}.) Present? (\emph{lǐwù}.) To give as a present? (\emph{sòng}.) A little while? (\textbf{\zh{一会儿}}, \emph{yíhuìr}.) See you tomorrow? (\textbf{\zh{明天见}}, \emph{míngtiān jiàn}.) To leave it till later? (\textbf{\zh{再说}}, \emph{zàishuō}.) Good night? (\emph{wǎn'ān}.) Bye-bye? (\emph{bàibai}.) Don't mention it? (\textbf{\zh{不客气}}, \emph{bú kèqi}.) To trouble somebody? (\emph{máfan}.) To thank? (\emph{gǎnxiè}.) Extremely? (\emph{fēicháng}.) Really? (\emph{zhēn}.) How old? (\textbf{\zh{多大}}, \emph{duō dà}.) To grow up? (\textbf{\zh{长大}}, \emph{zhǎngdà}.) Ago? (\textbf{\zh{前}}, \emph{qián}.) Age? (\emph{niánjì}.) In childhood? (\emph{xiǎoshíhou}.) To start class? (\textbf{\zh{上课}}, \emph{shàngkè}.) To finish class? (\textbf{\zh{下课}}, \emph{xiàkè}.) Classroom? (\emph{jiàoshì}.) Homework? (\emph{zuòyè}.) To take an exam? (\emph{kǎoshì}.)
\end{tcolorbox}
